\sect{Ćwiczenia 12}

\begin{frame}{Zadanie 1}
	\begin{figure}
		\includegraphics[scale=0.06, alt={}]{./imags/E12/01.png}
		\caption{Zadanie 1. Znajdź zastępcze źródło Thevenina i Nortona względem wyprowadzeń A i B. Zaprojektuj obwód stanowiący impedancję dopasowaną pozwalającą na pobranie maksymalnej mocy z obwodu.}
	\end{figure}	
\end{frame}

\begin{frame}{Zadanie 2}
	\begin{figure}
		\includegraphics[scale=0.06, alt={}]{./imags/E12/02.png}
		\caption{Zadanie 2. Znajdź zastępcze źródło Thevenina i Nortona względem wyprowadzeń A i B. Zaprojektuj obwód stanowiący impedancję dopasowaną pozwalającą na pobranie maksymalnej mocy z obwodu.}
	\end{figure}	
\end{frame}

\begin{frame}{Zadanie 3}
	\begin{figure}
		\includegraphics[scale=0.06, alt={}]{./imags/E12/03.png}
		\caption{Zadanie 3. Znajdź zastępcze źródło Thevenina i Nortona względem wyprowadzeń A i B. Zaprojektuj obwód stanowiący impedancję dopasowaną pozwalającą na pobranie maksymalnej mocy z obwodu.}
	\end{figure}	
\end{frame}

\begin{frame}{Zadanie 4}
	\begin{figure}
		\includegraphics[scale=0.06, alt={}]{./imags/E12/04.png}
		\caption{Zadanie 4. Znajdź zastępcze źródło Thevenina i Nortona względem wyprowadzeń A i B. Zaprojektuj obwód stanowiący impedancję dopasowaną pozwalającą na pobranie maksymalnej mocy z obwodu.}
	\end{figure}	
\end{frame}